\documentclass[table]{beamer}
%[]中可以使用handout、trancompress等参数

%指定beamer的模式与主题
\mode<presentation>
{
  \usetheme{Madrid}
%\usetheme{Boadilla}
%\usecolortheme{default}
%\usecolortheme{orchid}
%\usecolortheme{whale}
%\usefonttheme{professionalfonts}
}

%\usetheme{Madrid}
%这里还可以选择别的主题：Bergen, Boadilla, Madrid, AnnArbor, CambridgeUS, Pittsburgh, Rochester, Warsaw, ...
%有导航栏的Antibes, JuanLesPins, Montpellier, ...
%有内容的Berkeley, PaloAlto, Goettingen, Marburg, Hannover, ...
%有最小导航栏的Berlin, Ilmenau, Dresden, Darmstadt, Frankfurt, Singapore, Szeged, ...
%有章和节表单的Copenhagen, Luebeck, Malmoe, Warsaw, ...

%\usecolortheme{default}
%设置内部颜色主题(这些主题一般改变block里的颜色);这个主题一般选择动物来命名
%这里还可以选择别的颜色主题,如默认的和有特别目的的颜色主题default,structure,sidebartab,全颜色主题albatross,beetle,crane,dove,fly,seagull,wolverine,beaver

%\usecolortheme{orchid}
%设置外部颜色主题(这些主题一般改变title里的颜色);这个主题一般选择植物来命名
%这里还可以选择别的颜色主题,如默认的和有特别目的的颜色主题lily,orchid,rose

%\usecolortheme{whale}
%设置字体主题;这个主题一般选择海洋动物来命名
%这里还可以选择别的颜色主题,如默认的和有特别目的的颜色主题whale,seahorse,dolphin

%\usefonttheme{professionalfonts}
%类似的还可以定义structurebold,structuresmallcapsserif,professionalfonts


% 控制 beamer 的风格,可以根据自己的爱好修改
%\usepackage{beamerthemesplit} %使用 split 风格
%\usepackage{beamerthemeshadow} %使用 shadow 风格
%\usepackage[width=2cm,dark,tab]{beamerthemesidebar}


% 设定英文字体
\usepackage{fontspec}
%\setmainfont{Times New Roman}
%\setsansfont{Arial}
%\setmonofont{Courier New}

% 设定中文字体
\usepackage[BoldFont,SlantFont,CJKchecksingle,CJKnumber]{xeCJK}
%\setCJKmainfont[BoldFont={Adobe Heiti Std},ItalicFont={Adobe Kaiti Std}]{Adobe Song Std}
%\setCJKmainfont[BoldFont={Adobe Heiti Std},ItalicFont={Adobe Kaiti Std}]{WenQuanYi Micro Hei}
%\setCJKsansfont{Adobe Heiti Std}
%\setCJKmonofont{Adobe Fangsong Std}
%\punctstyle{hangmobanjiao}

\defaultfontfeatures{Mapping=tex-text}
\usepackage{xunicode}
\usepackage{xltxtra}

\XeTeXlinebreaklocale "zh"
\XeTeXlinebreakskip = 0pt plus 1pt minus 0.1pt

\usepackage{setspace}
\usepackage{colortbl,xcolor}
\usepackage{hyperref}
%\hypersetup{xetex,bookmarksnumbered=true,bookmarksopen=true,pdfborder=1,breaklinks,colorlinks,linkcolor=blue,filecolor=black,urlcolor=cyan,citecolor=green}
\hypersetup{xetex,bookmarksnumbered=true,bookmarksopen=true,pdfborder=1,breaklinks,colorlinks,linkcolor=cyan,filecolor=black,urlcolor=blue,citecolor=green}

% 插入图片
\usepackage{graphicx}
% 指定存储图片的路径（当前目录下的figures文件夹）
\graphicspath{{figures/}}

% 可能用到的包
\usepackage{amsmath,amssymb}
\usepackage{multimedia}
\usepackage{multicol}

% 定义一些自选的模板，包括背景、图标、导航条和页脚等，修改要慎重
% 设置背景渐变由10%的红变成10%的结构颜色
%\beamertemplateshadingbackground{red!10}{structure!10}
%\beamertemplatesolidbackgroundcolor{white!90!blue}
% 使所有隐藏的文本完全透明、动态，而且动态的范围很小
\beamertemplatetransparentcovereddynamic
% 使itemize环境中变成小球，这是一种视觉效果
\beamertemplateballitem
% 为所有已编号的部分设置一个章节目录，并且编号显示成小球
\beamertemplatenumberedballsectiontoc
% 将每一页的要素的要素名设成加粗字体
\beamertemplateboldpartpage

% item逐步显示时,使已经出现的item、正在显示的item、将要出现的item呈现不同颜色
\def\hilite<#1>{
 \temporal<#1>{\color{gray}}{\color{blue}}
    {\color{blue!25}}
}

% 自定义彩色块状结构的颜色
\setbeamercolor{bgcolor}{fg=yellow,bg=cyan}

% 在表格、图片等得标题中显示编号
\setbeamertemplate{caption}[numbered]

% 打开PDF后直接全屏
\hypersetup{pdfpagemode={FullScreen}}


% 使用 \part,\section,\subsection 等命令组织文档结构
% 使用 \frame 命令制作幻灯片

\begin{document}

\logo{\includegraphics[height=0.09\textwidth]{shu.png}}
\title[BN]{Batch Normalization}
\author[杜行健]{杜行健}
\institute[翔英学院 2015级]{翔英学院 2015级}
\date{\today}

% 定义目录页
\AtBeginPart{
  \frame{
    \frametitle{\partpage}
    \begin{multicols}{2}
% 如果目录过长，可以打开这个选项分两栏显示
      \tableofcontents
% 使用这个命令自动生成目录
    \end{multicols}
  }  
}


\begin{frame}
  \titlepage
\end{frame}

\begin{frame}[plain]
        \frametitle{Overview}
  \setcounter{tocdepth}{2}
  \tableofcontents
\end{frame}

\section{Issues With Training Deep Neural Networks}
\begin{frame}
    \frametitle{Covariate Shift}
        Given the same observation $X=x$, the conditional distributions of $Y$ are the same in the source and target domains. However, the marginal distributuions of X may be different in the source and the target domains. Formally, we assume that $P_s (Y|X=x) = P_t (Y|X = x) for all x \in \chi$, but $P_s (X) \neq P$. This difference between the two domains is called \textit{covariate shift}.\par
\end{frame}
\begin{frame}
        \frametitle{Internal Covariate Shift}
        In the case of deep networks, the input to each layer is affected by parameters in all the input layers. So even small changes to the network get amplified down the network. This leads to change in the input distribution to internal layers of the deep network and is known as internal covariate shift.\par
        It is well established that networks converge faster if the inputs have been whitened (ie zero mean, unit variances) and are uncorrelated and internal covariate shift leads to just the opposite.\par
\end{frame}

\begin{frame}
        \frametitle{Vanishing Gradient}
        Saturating nonlinearities (like tanh or sigmoid) can not be used for deep networks as they tend to get stuck in the saturation region as the network grows deeper. Some ways around this are to use:
        \begin{itemize}
                \hilite <1> \item Nonlinearities like ReLU which do not saturate
                \hilite <2> \item Smaller learning rates
                \hilite <3> \item Careful initializations

        \end{itemize}
\end{frame}


\section{Batch Normalization}
\begin{frame}
    \frametitle{Standard Score}
        Let us say that the layer we want to normalize has $d$ dimensions $X = (x_1, \cdot, x_d)$. Then, we can normalize the $k^{th}$ dimension as follows:\par
        \begin{equation}
                \hat{x}^{k} = \frac{x^k - E[x^k]}{\sqrt{Var[x^k]}}
        \end{equation}
\end{frame}

\begin{frame}
        We also need to scale and shift the normalized values otherwise just normalizing a layer would limit the layer in terms of what it can represent. For example, if we normalize the inputs to a sigmoid function, then the output would be bound to the linear region only. \par
        So the normalized input $x^k$ is transformed to:
        \begin{equation}
                y^k = \lambda^{k} \hat{x} + \beta^{k}
        \end{equation}
        where $\lambda$ and $\beta$ are parameters to be learned.\par
\end{frame}

\begin{frame}
        Moreover, just like we use mini-batch in Stochastic Gradient Descent (SGD), we can use mini-batch with normalization to estimate the mean and variance for each activation.\par
        The transformation from x to y as described above is called Batch Normalizing Tranform. This BN transform is differentiable and ensures that as the model is training, the layers can learn on the input distributions that exhibit less internal covariate shift and can hence accelerate the training.\par
        At training time, a subset of activations in specified and BN transform is applied to all of them.\par
        During test time, the normalization is done using the population statistics instead of mini-batch statistics to ensure that the output deterministically depends on the input.\par
\end{frame}
\section{Advantages Of Batch Normalization}
\begin{frame}
        \frametitle{Advantages Of Batch Normalization}
		 \begin{itemize}
                \hilite <1> \item Reduces internal covariant shift.
                \hilite <2> \item Reduces the dependence of gradients on the scale of the parameters or their initial values.
                \hilite <3> \item Regularizes the model and reduces the need for dropout, photometric distortions, local response normalization and other regularization techniques.
				\hilite <4> \item Allows use of saturating nonlinearities and higher learning rates.

        \end{itemize}

\end{frame}
\begin{frame}

	\begin{center}
	\Huge{\textbf{Thanks for your attention!}}

	\Huge{\textit{Any questions?}}
\end{center}
\end{frame}

\end{document}
